\chapter{Introduction générale}
\label{chap:introduction}
\justifying
\sloppy

Dans de nombreux systèmes de détection, le premier objectif consiste simplement à savoir qu'un événement a eu lieu. Cette information binaire est déjà utile, mais elle devient rapidement insuffisante dès que l'on cherche à reconstruire plus finement ce qui s'est passé. Il faut alors connaître non seulement la présence d'un signal, mais aussi son instant d'apparition, sa position et sa relation avec les autres signaux observés.

Cette idée est particulièrement importante pour les détecteurs rapides. Une impulsion peut être très courte, arriver à un instant imprévisible et devoir être associée à d'autres informations avant d'être exploitable. L'électronique de lecture ne peut donc pas se contenter de compter les événements : elle doit aussi les dater, les organiser et les transmettre dans un format robuste.

Dans ce cadre s'inscrit le projet \textbf{SPADMIC}. Il s'agit d'un prototype d'\textbf{ASIC} (\textit{Application-Specific Integrated Circuit}), c'est-à-dire un circuit intégré conçu pour une application dédiée. Il vise à intégrer une électronique de lecture numérique autour d'une matrice de \textbf{SPAD} (\textit{Single-Photon Avalanche Diode}), des détecteurs capables de produire une impulsion lorsqu'un photon est détecté. Dans ce type de système, la mesure temporelle devient une information centrale.

Pour mesurer un intervalle de temps dans une chaîne numérique, on utilise un convertisseur temps-numérique, ou \textbf{TDC} (\textit{Time-to-Digital Converter}). Le bloc étudié est une version multiphase de ce type de convertisseur, appelée \textbf{MPTDC} (\textit{Multi-Phase Time-to-Digital Converter}). La progression retenue construit d'abord le contexte : rôle de la puce \textbf{SPADMIC}, besoins de mesure temporelle, principes des \textbf{TDC}, puis implémentation RTL détaillée du \textbf{MPTDC}. Le terme \textbf{RTL} (\textit{Register Transfer Level}) désigne ici la description numérique du circuit avant synthèse matérielle.

La figure \ref{fig:spadmic-mind-model} donne un premier modèle mental de cette organisation. Elle situe la matrice \textbf{SPAD}, les trois voies \textbf{MPTDC}, le bloc de position, l'arbitrage numérique et la calibration hôte avant d'entrer dans les détails techniques des chapitres suivants.

\begin{figure}[H]
    \centering
  \includegraphics[width=0.92\textwidth,trim=0 105 0 0,clip]{spadmic/spadmic_mind_model.png}
    \caption{Vue d'ensemble du prototype \textbf{SPADMIC} et de la place du cœur \textbf{MPTDC}.}
    \label{fig:spadmic-mind-model}
\end{figure}

\section{Contexte et problématique}

La conception d’un \textbf{TDC} de haute précision impose de concilier plusieurs exigences souvent difficiles à satisfaire simultanément. Il faut réduire le pas de quantification, conserver une plage dynamique utile, limiter le bruit et le jitter, et maintenir une architecture compatible avec une intégration dans un \textbf{ASIC}. Dans ce cadre, les architectures de type Vernier occupent une place particulièrement intéressante, car elles permettent de s’affranchir partiellement de la limite imposée par le délai élémentaire d’une cellule logique en exploitant la différence entre deux délais voisins.

Dans \textbf{SPADMIC}, la présence de trois voies \textbf{TDC} ne vient pas d'un simple besoin de paralléliser trois événements indépendants. Elle est liée à l'organisation tri-axis déjà validée dans le cadre de \textbf{PICMIC} (\textit{PICosecond-MICrometer sensor})~\cite{abreu2024picmic0}. Ce concept repose sur une matrice de pixels hexagonaux interconnectés : la charge ou l'impulsion générée par un impact est projetée sur trois familles de lignes orientées selon \SI{-120}{\degree}, \SI{0}{\degree} et \SI{120}{\degree}. Cette lecture projective réduit fortement le nombre de canaux électroniques actifs par rapport à une lecture pixel par pixel, tout en conservant assez d'information géométrique pour reconstruire la position par croisement des projections et calcul de centroïde. Les trois \textbf{TDC} correspondent donc aux trois axes de lecture d'un même événement physique, avant mutualisation de la sortie numérique.

La thèse d’Amina Annagrebah, utilisée ici comme référence théorique principale, propose une étude approfondie de ces architectures et, plus précisément, d’une approche Vernier à oscillateurs en anneau multiphase appliquée à des besoins de mesure temporelle de très haute précision \cite{annagrebah2021thesis}. Ce travail constitue une base particulièrement solide pour introduire les notions nécessaires avant d’aborder l’architecture développée dans le cadre du projet \textbf{SPADMIC} et son cœur \textbf{MPTDC}.

\section{Périmètre et organisation du document}

Le périmètre retenu reste volontairement maîtrisé. Au moment de la rédaction initiale, l’architecture globale de \textbf{SPADMIC} n’était pas encore totalement figée, même si elle s’est depuis davantage stabilisée. Il ne s’agit donc pas de décrire tout l’ASIC comme un objet déjà définitivement arrêté, mais de présenter de manière cohérente ce qui est déjà structuré, justifiable et techniquement défendable.

La progression commence par le cadre scientifique et institutionnel au sein duquel s’inscrit l’alternance. Elle introduit ensuite le contexte général du projet \textbf{SPADMIC}, les contributions numériques réalisées et la place du \textbf{MPTDC} dans la chaîne de lecture de l’ASIC. La partie suivante développe les fondements théoriques des \textbf{TDC}, leurs métriques de performance et les architectures de type Vernier. Sur cette base, l’étude se prolonge par la description de l’architecture de référence extraite de la thèse, puis par une transition vers l’implémentation effectivement retenue. Enfin, l'analyse détaille l’architecture active de \textbf{MPTDC}, puis présente la méthodologie de caractérisation, les résultats de calibration et les limites encore ouvertes vers une future intégration plus aboutie.

\clearpage
